\documentclass[book.tex]{subfiles}

\begin{document}

\chapter{Units, Types and Dimensional Analysis}
\label{chap:units_and_types}
\noindent\rule{11cm}{0.4pt} \\
\textbf{Prerequisites:} \autoref{chap:simple_physics}, \autoref{chap:type_theory} \\
\textbf{Difficulty Level:} * \\
\noindent\rule{11cm}{0.4pt}
\vspace{5mm}

One of the most useful applications of type theory is to track units and dimensional analysis for physical systems. In this chapter, we will explain how to use a type system to track units and perform dimensional analysis with Physika.

We start by considering mass. We say that mass $m$ is in kilograms by assigning it the type
\begin{lstlisting}
m: kg
\end{lstlisting}
Here \lstinline{kg} is a modified version of type $\mathbb{R}^+$ with a modified unit (the unit corresponds to one kilogram). Let's next consider length. We say length \lstinline{l} is denoted in meters by assigning type
\begin{lstlisting}
l: meter
\end{lstlisting}
Here \lstinline{m} is a modified version of $\mathbb{R}^+$ with a modified unit (the unit corresponds to one meter). Similarly, we say \lstinline{t} is a time interval measured in seconds by saying
\begin{lstlisting}
t: seconds
\end{lstlisting}
Units form a mathematical algebra. That is, the following quantities are all units
\begin{lstlisting}
m * s, m / s, m / s^2, kg * m / s^2
\end{lstlisting}
Physika's type system allows us to track the units of expressions.

An alternative method for facilitating dimensional analysis is to have units have more basic types

\begin{lstlisting}
l: [L]
t: [T]
m: [M]
\end{lstlisting}

Here we use \lstinline{[L],[T], [M]} to denote the most basic types for distance, time and mass rather than a specific instantiation such as meters, seconds, or kilograms. This more abstract level of analysis can prove useful when doing abstract analyses of systems.

An important concept that arises is that units can cancel. We can model this by stating that for example

\begin{lstlisting}
a : meter / meter $\sim a$ : $\mathbb{R}$
\end{lstlisting}
That is, a quantity that is of type \lstinline{meter/meter} is just a scalar quantity. We can then proceed to write down types for common physical quantities

\begin{lstlisting}
Area: [L]$^2$
Speed: [L][T]$^{-1}$
Acceleration: [L][T]$^{-2}$
Force: [M][L][T]$^{-2}$
Energy: [M][L]$^2$[T]$^{-2}$
\end{lstlisting}

This type of decomposition can be very useful when studying the behavior of physical systems. For example, we can derive the types of physical constants from this analysis. Consider Newton's gravitational equation
\begin{align}
    F &= - \frac{G m_1 m_2}{r^2}
\end{align}
We can consider the types on both sides to obtain the equation
\begin{lstlisting}
[M][L][T]$^{-2}$ = [G][M]$^2$[L]$^{-2}$
\end{lstlisting}
Solving, we obtain
\begin{lstlisting}
[G] = [M]$^{-1}$[L]$^3$[T]$^{-2}$
\end{lstlisting}
We can solve similarly for the types of other physical constants.
\section{Buckingham $\pi$ Theorem}

Suppose that we have a physical equation given by

\begin{align}
    f(q_1,\dotsc,q_n) &= 0
\end{align}

Assume that the $q_i$ are physically independent variables and assume there are $k$ units (mass, time, etc) in this equation. Then the Buckingham-$\pi$ theorem states that there are $p = n - k$ dimensionless terms
\begin{align}
    \pi_i &= q_1^{a_1}\dotsc q_n^{a_n}
\end{align}
where the $a_i$ are rational terms such that the original equation can be restated as 

\begin{align}
    F(\pi_1,\dotsc,\pi_p) &= 0
\end{align}
Philosophically, the Buckingham-$\pi$ theorem demonstrates that the laws of physics are independent of particular choices of units. 

%\textcolor{red}{TODO: Add a proof of the Buckingham $\pi$ theorem and relationship with types.}
\end{document}